% Conclusion. Source: HTML §9; synthesis only (paper-writing.mdc §6).

\section{Conclusion}
\label{sec:conclusion}

This report addressed target-speaker recognition in recordings where the enrolled speaker may or may not talk.
We proposed \tsbench, which will be released publicly.
Every condition is paired with target-present and target-absent samples.
The synthetic part is designed from errors observed in real use and changes one factor at a time, and the real part comes from the test sets of public meeting and call corpora and from web conversations.
With \tsbench, we showed that evaluations limited to a single talker, to samples with the target present or to silence can rate two systems as roughly equal.
Yet when others talk in real recordings, one of those systems almost always produces false alarms and the other rarely does.
We also presented AmphionASR~2.0, which joins the enrollment and the mixture with a learnable separator after the convolutional downsampling of its audio encoder and whose training data give every scene type target-absent counterparts.
Compared on \tsbench with Xiaomi-CocktailASR-1, a model about five times larger, AmphionASR~2.0 has a lower CER with the target present and a false-alarm rate more than an order of magnitude lower when others talk and the target is absent.
The difference holds on real meetings and calls.
The same weights also handle plain and hotword ASR, with a small regression in general recognition relative to the initial model.
False alarms when a similar-sounding person speaks briefly, near-homophone distractors in hotword lists, and this regression remain the main open problems.
